\section{本章小结}
\label{sec:func-summary}

章首的问题是：未知量是一整条曲线时，怎样找最好的那一条。函数可以看成很长的向量，泛函梯度对应向量梯度，这个类比决定了本章的工作顺序：先证存在，再用微分判别，最后离散计算。

\paragraph{三件工具各负责什么}
弱极限与下半连续性负责证明存在。无限维的闭球不紧，只能取弱收敛子列，于是可行集要弱序列闭，目标要弱序列下半连续。凸性让这两个条件可以检查，定理~\ref{thm:func-direct-method}把它们串起来。第~\ref{sec:func-existence}节的四个反例分别对应缺少强制性、反射性、弱闭性和对导数的凸性。写出驻定方程不能代替这一步。

微分与变分负责判别候选。方向导数、Gâteaux微分和Fréchet微分逐层加强，差别在于余项是否对方向一致，并且依赖范数。微分经Riesz表示才成为梯度，梯度依赖内积，也可能不存在。一阶条件和凸性充分性与有限维一致，二阶充分性需要统一的曲率下界，乘子法则在有限个约束下成立，无穷约束下要重新检查拓扑。

试探空间负责把函数问题转成有限维计算。Ritz方法在二次目标上给出能量范数下的最佳逼近。固定函数的能量逼近、离散极小者的收敛、数值求解的误差是三件不同的事。

\paragraph{积分目标的结论}
积分泛函的一阶变分经分部积分，分成内部的Euler--Lagrange方程和一个边界项。内部方程来自两端固定的扰动，边界条件来自放开的扰动：端点值自由给出$F_p=0$，端点沿曲线移动给出横截条件。这些方程都是必要条件，极小性要靠凸性或二阶条件来确认。

\begin{BookInlineTable}
  \centering
  \caption{本章的问题、工具与条件。}
  \label{tab:func-summary}
  \begin{tabular}{lll}
    \toprule
    问题 & 工具 & 关键条件\\
    \midrule
    解是否存在 & 直接法 & 反射、弱闭、弱下半连续\\
    候选怎样判别 & 微分与变分 & 所选范数下可微、凸性\\
    方程与边界 & 分部积分 & 正则性、允许的扰动\\
    怎样计算 & Ritz方法 & 稠密、连续、求解精度\\
    \bottomrule
  \end{tabular}
\end{BookInlineTable}

\paragraph{边界}
本章只建立优化问题本身的存在性、判别和近似。Euler--Lagrange方程解的存在性、正则性和具体边值问题的理论属于微分方程部分；以概率分布为未知量的变分问题由\BookChapterRef{chap:probability-distribution-optimization}{卷七第12章“概率分布的优化与逼近”}研究。离散之后的有限维问题，用本卷前面各部分的方法求解。
